Bài báo cáo trình bày quy trình phân tích và tiền xử lý dữ liệu FoodSeg103 cho bài toán phân vùng đa lớp, đồng thời thiết lập baseline BiSeNet trên bộ dữ liệu đã cân bằng lại. Qua thống kê cho thấy dữ liệu gốc có chất lượng cao nhưng bị long-tail nghiêm trọng, dẫn đến nguy cơ thiên lệch khi huấn luyện. Để giảm tác động này, chúng tôi áp dụng gộp nhãn theo ngữ nghĩa và lọc nhiễu để tái cấu trúc không gian nhãn từ 103 xuống 76 lớp, đồng thời giữ được phần lớn mẫu huấn luyện. Dựa trên bộ dữ liệu sau xử lý, pipeline baseline BiSeNet được chuẩn hóa với các bước nạp dữ liệu, augmentation, tối ưu và đánh giá định kỳ theo mIoU/mAcc/pixel accuracy. Cấu trúc này tạo mốc tham chiếu rõ ràng cho các hướng cải tiến tiếp theo hướng tới cân bằng giữa độ chính xác và tốc độ suy luận.